\section{Acknowledgments}

The research was supported by the Russian Science Foundation grant No. 25-11-00355, \url{https://rscf.ru/project/25-11-00355/}. 

\section{Conclusion and discussion}

This paper introduces a complexity-aware fine-tuning pipeline that measures model response uncertainty using the entropy of its predicted answer, then trains on the resulting easy, medium, and hard splits via different tactics.

We confirm that entropy works as a difficulty estimation. Single-token answer entropy reaches ROC AUC values up to $0.8$, clearly beating MASJ-based estimates of $0.57$. 
This confirms that a model's own confidence is a reliable, automatic proxy for question difficulty.

Using the entropy-based data split, we find that different complexity scores require different training approaches. Standard supervised fine-tuning (SFT) is enough for the easy and medium bands, but lags on the hard band. For hard questions, adding a distilled chain of thought from a large LLM unlocks further gains. Our pipeline achieves accuracies of 0.52/0.64 vs. 0.39/0.51 for Qwen 3B/Phi-4-mini, while using $81\%$ less data.

The pipeline is fully automated and can be included in other fine-tuning workflows. Our findings suggest that curriculum ideas still matter for today's LLMs: letting the model focus on what it can already solve directly, while giving extra guidance only where it struggles, yields a better allocation of limited model capacity.



